%! Graphing Quadratic Functions & Transformations — Unit 2, Lesson 2.1 — Guided Notes & Practice
% THE in-class document, in the MAIN IDEAS / NOTES shape (LESSON_SHAPE.md §1): the vocabulary
% box, the hook, then ONE two-column guidednotes table. Each row is one idea — a label on the
% left; on the right one or two COMPLETE printed sentences, ONE large pre-drawn display, then a
% two-across grid of numbered problems with real writing room. Problems run 1..14 through the
% component. No group activity: the notes carry the whole gradual release, closed by a SPOKEN
% debrief; the homework is the individual practice.
%   I DO   : the teacher reads each row's definition, marks up its display, and works the FIRST
%            problem of each grid (1, 3, 6, 9).
%   WE DO  : the class works the rest of each grid, students holding the pen; the last row,
%            Guided Practice (11–14), is worked entirely together on a new curve with no story.
%   YOU DO : nothing on this page — ../homework/ IS the individual practice, started in the
%            last ten minutes of the period.
% DENSITY: no mid-sentence blanks — the only \blank{}s are the two hook questions, the cells of
% the value table in problem 3 with its one-word follow-up, and problem 8's rule-and-intercept.
% Every display is pre-drawn and read or labelled (\labelbox); THERE ARE NO SKETCH QUESTIONS.
% ONE CONTEXT: the unit anchor from 2.0 — x^2-2x-3 = (x-1)^2-4 = (x+1)(x-3), one curve in three
% costumes. Vertex (1,-4), axis x=1, zeros -1 and 3, y-intercept (0,-3).
% THE CRUX is the HORIZONTAL-SHIFT SIGN FLIP: (x-1)^2 moves the parent RIGHT, while in the
% factored form (x+1)(x-3) the zeros are -1 and 3 — the same bracket does opposite things in the
% two forms, and which form you are in decides. Row 4 settles it (problems 9–10); Guided
% Practice problem 14 tests it again on new ground.
% NOTHING IS FACTORED OR SOLVED HERE — 2.2 factors, 2.3–2.5 solve. A factored form is GIVEN.
% PAGE LOCKSTEP with notes_key/: \vterm -> \vtermans, \blank -> \ans, and the answer argument of
% every \pcell and \labelbox filled; every \begin{work} block is BYTE-IDENTICAL in both files.
% Nothing else differs. TABLE MECHANICS: inside a cell `\\` ENDS THE ROW — break lines with
% \par; the figure and EACH grid row sit in their own sub-row (`\\` then `& ...`, the grid
% reopening as probgrid*) so the table can break between them.
\documentclass[12pt]{article}
\usepackage{algebra2-article}
\usepackage{algebra2-boxes}
% Vocabulary rows (SAAR style, user correction 2026-09-06): a FIXED-HEIGHT row carrying the term
% label and OPEN writing space — no inline blank and no rule line. The key fills exactly the same
% height (definitions two lines at most), so the box cannot drift blank vs. keyed. Six rows at
% 1.5cm, so the box and the hook still share page 1.
\newcommand{\vrowheight}{2.0cm}
\newcommand{\vterm}[1]{%
  \par\noindent\begin{minipage}[t][\vrowheight][t]{\linewidth}%
    \vspace*{5pt}\textbf{\textcolor{forest}{#1:}}%
  \end{minipage}\par}
% Coordinate grid: {xmin}{xmax}{ymin}{ymax}
\newcommand{\qgrid}[4]{%
  \draw[step=1, linegray!40, very thin] (#1,#3) grid (#2,#4);
  \draw[->, semithick, charcoal] (#1,0)--(#2,0) node[below right, font=\scriptsize]{$x$};
  \draw[->, semithick, charcoal] (0,#3)--(0,#4) node[left, font=\scriptsize]{$y$};}
% Points: {color}{x}{y}
\newcommand{\fdot}[3]{\fill[#1] (#2,#3) circle (3.2pt);}
\newcommand{\hdot}[3]{\draw[very thick, #1, fill=white] (#2,#3) circle (3.2pt);}

\begin{document}

\pageheader{Unit 2, Lesson 2.1}{Guided Notes \& Practice}
% NAMESTRIP: no \namedateperiod here — the name row lives on the cover only.

\vspace{6pt}

% ── PAGE 1: vocabulary, then the hook ────────────────────────────────────────
\begin{vocabbox}
\small
Fill in each term \textbf{as we name it} in the notes below.
\vterm{Parent function $y=x^2$}
\vterm{Transformation (shift, stretch, reflection)}
\vterm{Vertex form $a(x-h)^2+k$}
\vterm{Standard form $ax^2+bx+c$}
\vterm{Intercept (factored) form $a(x-p)(x-q)$}
\end{vocabbox}

\vspace{8pt}

\boxguard[18]
\begin{hookbox}
\small
\begin{minipage}[c]{0.52\linewidth}
\textbf{One curve in three costumes.} In Lesson~2.0 you \emph{read} this parabola. Here are
three expressions --- and \emph{one} graph.
\[ x^2-2x-3 \;=\; (x-1)^2-4 \;=\; (x+1)(x-3) \]
\end{minipage}\hfill
\begin{minipage}[c]{0.44\linewidth}\centering
\begin{tikzpicture}[scale=0.32]
  \qgrid{-3.5}{5.5}{-5.5}{5.5}
  \begin{scope}
    \clip (-3.5,-5.5) rectangle (5.5,5.5);
    \draw[very thick, forest, domain=-2.4:4.4, samples=80, smooth, variable=\x]
      plot (\x,{(\x-1)^2-4});
  \end{scope}
  \fdot{forest}{1}{-4}
\end{tikzpicture}
\end{minipage}
\par\vspace{5pt}
\textbf{Three equations. Circle one:} \quad \textbf{three different parabolas} \qquad
\textbf{one parabola wearing three costumes}
\par\vspace{4pt}
Which costume tells you the \textbf{vertex} fastest? \blank{4.0cm} \quad Which one tells you the
\textbf{zeros}? \blank{4.0cm}
\par\vspace{4pt}
Every question today is this picture: \textbf{whichever form is in front of you, what does it
hand you for free --- and how has the parent $y=x^2$ been moved to get here?}
\end{hookbox}

\small
% Displays inside table cells: tight skips, or every brace costs a centimetre. (\small resets
% the display skips, so this must follow it.) The work blocks use the same skips. Figures are
% centred with \centering inside the cell, not a center environment (its \topsep is ~1cm a figure).
\setlength{\abovedisplayskip}{4pt}\setlength{\belowdisplayskip}{4pt}%
\setlength{\abovedisplayshortskip}{2pt}\setlength{\belowdisplayshortskip}{2pt}%
\begin{guidednotes}
% ── ROW 1: the parent and its shifts — the sign flip is planted here ─────────
\mainidea[The parent and]{Its Shifts} &
The \textbf{parent function} of every parabola is $y=x^2$. A \textbf{transformation} changes the
parent's rule and so moves its graph. Adding $k$ \emph{outside}, $f(x)+k$, slides it \textbf{up
or down}; changing $x$ to $x-h$ \emph{inside}, $f(x-h)$, slides it \textbf{left or right} ---
and \textbf{subtracting inside moves it \emph{right}}, the opposite of what the sign looks like.
\\
& {\centering
\begin{tikzpicture}[scale=0.22]
  \qgrid{-3.5}{5.5}{-1.5}{8.5}
  \begin{scope}
    \clip (-3.5,-1.5) rectangle (5.5,8.5);
    \draw[very thick, charcoal, domain=-3.2:3.2, samples=80, smooth, variable=\x]
      plot (\x,{(\x)^2});
    \draw[very thick, forest, domain=-2.6:2.6, samples=80, smooth, variable=\x]
      plot (\x,{(\x)^2+3});
    \draw[very thick, redacc, domain=-1.2:5.2, samples=80, smooth, variable=\x]
      plot (\x,{(\x-2)^2});
  \end{scope}
  \fdot{charcoal}{0}{0} \fdot{forest}{0}{3} \fdot{redacc}{2}{0}
  \node[left, font=\scriptsize\bfseries, charcoal] at (-2.6,7.0) {$y=x^2$};
  \node[above, font=\scriptsize\bfseries, forest] at (0,3.4) {$y=x^2+3$};
  \node[below right, font=\scriptsize\bfseries, redacc] at (2.2,-0.2) {$y=(x-2)^2$};
\end{tikzpicture}\par\vspace{2pt}
$(x-2)^2$ moves the parent \labelbox{2.6cm}{}\par}
\\
& \notesprompt{Outside moves it up and down; inside moves it left and right.}
\begin{probgrid}{|Y|Y|}
\pcell{1}{$y=x^2+3$ --- what does the $+3$ do to the parent, and where is the vertex
now?}{1.8cm}{} &
\pcell{2}{$y=(x-2)^2$ --- which way does it move, and how far? Why does the sign look
wrong?}{1.9cm}{} \\ \hline
\end{probgrid}
\\ \hline
% ── ROW 2: a — stretch, squash, flip ─────────────────────────────────────────
\mainidea[Stretch and]{Flip} &
The number $a$ in $a\,f(x)$ multiplies every output. If $|a|>1$ the parabola is
\textbf{narrower} than the parent; if $0<|a|<1$ it is \textbf{wider}; and if $a<0$ it is
\textbf{reflected} --- it opens \emph{down}, so its vertex is a \textbf{maximum} instead of a
minimum (Lesson~2.0's language, now read straight off the equation).
\\
& {\centering
\begin{tikzpicture}[scale=0.22]
  \qgrid{-4.5}{4.5}{-4.5}{8.5}
  \begin{scope}
    \clip (-4.5,-4.5) rectangle (4.5,8.5);
    \draw[very thick, charcoal, domain=-3.2:3.2, samples=80, smooth, variable=\x]
      plot (\x,{(\x)^2});
    \draw[very thick, forest, domain=-1.9:1.9, samples=80, smooth, variable=\x]
      plot (\x,{3*(\x)^2});
    \draw[very thick, redacc, domain=-3.2:3.2, samples=80, smooth, variable=\x]
      plot (\x,{-0.5*(\x)^2});
  \end{scope}
  \fdot{charcoal}{0}{0}
  \node[right, font=\scriptsize\bfseries, charcoal] at (2.9,8.0) {$y=x^2$};
  \node[right, font=\scriptsize\bfseries, forest] at (1.7,7.4) {$y=3x^2$};
  \node[right, font=\scriptsize\bfseries, redacc] at (3.1,-4.0) {$y=-\tfrac12 x^2$};
\end{tikzpicture}\par\vspace{2pt}
All three share the same vertex; only $a$ is different.\par}
\\
& \textbf{3.} Check the word ``narrower'' with numbers, not with your eyes --- fill the table.
Then: at $x=2$, which graph is \emph{higher}? \blank{2.0cm}
\par\vspace{2pt}
\renewcommand{\arraystretch}{1.2}
\begin{tabularx}{\linewidth}{@{} l Y Y Y Y @{}}
\rowcolor{sky}
\textbf{$x$} & $0$ & $1$ & $2$ & $3$ \\
\hline
\textbf{$y=x^2$} & \blank{0.9cm} & \blank{0.9cm} & \blank{0.9cm} & \blank{0.9cm} \\
\textbf{$y=3x^2$} & \blank{0.9cm} & \blank{0.9cm} & \blank{0.9cm} & \blank{0.9cm} \\
\end{tabularx}
\\
& \notesprompt{Read $a$, then say direction and width.}
\begin{probgrid}{|Y|Y|}
\pcell{4}{$y=-\tfrac12 x^2$ --- opens which way, narrower or wider than the parent, and is its
vertex a maximum or a minimum?}{1.9cm}{} &
\pcell{5}{$y=4x^2$ and $y=\tfrac14 x^2$ share a vertex. Which is narrower, and what in $a$ told
you?}{1.9cm}{} \\ \hline
\end{probgrid}
\\ \hline
% ── ROW 3: vertex form reads itself ─────────────────────────────────────────
\mainidea[Vertex form]{Reads Itself} &
\textbf{Vertex form} is $y=a(x-h)^2+k$. It hands you the \textbf{vertex} $(h,k)$ and the
\textbf{axis of symmetry} $x=h$ with no work at all, and $a$ still gives direction and width.
\\
& {\centering
\begin{tikzpicture}[scale=0.22]
  \qgrid{-3.5}{5.5}{-5.5}{5.5}
  \draw[dashed, very thick, redacc] (1,-5.3)--(1,5.3);
  \begin{scope}
    \clip (-3.5,-5.5) rectangle (5.5,5.5);
    \draw[very thick, forest, domain=-2.4:4.4, samples=80, smooth, variable=\x]
      plot (\x,{(\x-1)^2-4});
  \end{scope}
  \fdot{forest}{1}{-4}
  \node[right, font=\scriptsize\bfseries, forest] at (1.3,-4.2) {vertex $(1,-4)$};
  \node[above right, font=\scriptsize\bfseries, redacc] at (1.1,3.4) {$x=1$};
\end{tikzpicture}\par\vspace{2pt}
The dashed line is the \labelbox{3.4cm}{}\par}
\\
& \textbf{Reading $y=a(x-h)^2+k$:} \quad
\stepnum{1}\ Read $a$ --- direction and width. \quad
\stepnum{2}\ Read $h$ --- \textbf{flip the sign inside}. \quad
\stepnum{3}\ Read $k$ --- no flip. \quad
\stepnum{4}\ Plot $(h,k)$, draw the axis $x=h$. \quad
\stepnum{5}\ Use symmetry for a point either side.
\\
& \notesprompt{Read the vertex straight off the form.}
\begin{probgrid}{|Y|Y|}
\pcell{6}{$y=(x-1)^2-4$ --- name $a$, $h$, and $k$. Vertex? Axis? Opens which way?}{1.8cm}{} &
\pcell{7}{$y=-2(x+3)^2+5$ --- vertex, axis, direction, and narrower or wider than the
parent.}{1.9cm}{} \\ \hline
\end{probgrid}
\\
& \textbf{8.} Run it backwards: a parabola has vertex $(4,-1)$, opens up, and is the same width
as the parent. Its rule in vertex form is $y=$ \blank{3.2cm}, and its $y$-intercept is
\blank{1.8cm}.
\\ \hline
% ── ROW 4: THE CRUX — three costumes, one curve, and the bracket that lies ───
\mainidea[Three costumes,]{One Curve} &
\textbf{Standard form} $y=ax^2+bx+c$ gives the $y$-intercept $(0,c)$ free; its axis is
$x=-\frac{b}{2a}$, and that $x$ substituted back gives the vertex. \textbf{Intercept (factored)
form} $y=a(x-p)(x-q)$ gives the \textbf{zeros} $p$ and $q$, with the axis at their
\emph{midpoint} $x=\frac{p+q}{2}$. \textbf{All three forms of our anchor agree --- it is one
curve.}
\\
& {\centering
\begin{tikzpicture}[scale=0.22]
  \qgrid{-3.5}{5.5}{-5.5}{5.5}
  \draw[dashed, very thick, redacc] (1,-5.3)--(1,5.3);
  \begin{scope}
    \clip (-3.5,-5.5) rectangle (5.5,5.5);
    \draw[very thick, forest, domain=-2.4:4.4, samples=80, smooth, variable=\x]
      plot (\x,{(\x-1)^2-4});
  \end{scope}
  \fdot{forest}{1}{-4} \fdot{navy}{-1}{0} \fdot{navy}{3}{0} \fdot{goldacc}{0}{-3}
  \node[above left, font=\tiny, navy] at (-1,0.1) {$(-1,0)$};
  \node[above right, font=\tiny, navy] at (3,0.1) {$(3,0)$};
  \node[left, font=\tiny, goldacc] at (-0.15,-3) {$(0,-3)$};
\end{tikzpicture}\par\vspace{2pt}
$x^2-2x-3\to(0,-3)$; \; $(x+1)(x-3)\to-1$ and $3$; \; $(x-1)^2-4\to(1,-4)$.\par}
\\
& \textbf{The axis, three ways.} From standard form; then check it against the midpoint of the
zeros, $\frac{-1+3}{2}=1$, and against $h=1$ in vertex form:
\begin{work}
  x &= -\frac{b}{2a} = -\frac{-2}{2(1)} = 1 \\
  y &= (1)^2 - 2(1) - 3 = -4
\end{work}
\textbf{Same vertex $(1,-4)$, same axis $x=1$, every time. But watch the bracket:} in
\emph{vertex} form $(x+4)^2$ shifts \emph{left}~$4$, while as a \emph{factor} $(x+4)$ gives the
zero $x=-4$ --- opposite-looking jobs, and \textbf{which form you are in decides.}
\\
& \notesprompt{Settle the sign flip.}
\begin{probgrid}{|Y|Y|}
\pcell{9}{A classmate says ``$y=(x+4)^2$ shifts the parent \emph{right}~$4$.'' Which way does
it really go, and why?}{1.9cm}{} &
\pcell{10}{Same bracket, two forms. In $y=(x+4)^2$ the shift is which way? In $y=(x+4)(x-2)$
the zeros are what? Why is that not the same question?}{1.9cm}{} \\ \hline
\end{probgrid}
\\ \hline
% ── ROW 5: guided practice — WE DO, together, a new curve in three costumes ──
\mainidea[Guided practice]{One More Curve} &
\begin{minipage}[c]{0.56\linewidth}
A new curve, no story --- and again in three costumes.
\par\vspace{2pt}
$y=-(x-2)^2+9=-x^2+4x+5=-(x+1)(x-5)$
\end{minipage}\hfill
\begin{minipage}[c]{0.40\linewidth}\centering
\begin{tikzpicture}[scale=0.19]
  \qgrid{-2.5}{6.5}{-2.5}{10.5}
  \begin{scope}
    \clip (-2.5,-2.5) rectangle (6.5,10.5);
    \draw[very thick, forest, domain=-1.5:5.5, samples=80, smooth, variable=\x]
      plot (\x,{-((\x-2)^2)+9});
  \end{scope}
  \fdot{forest}{2}{9} \fdot{navy}{-1}{0} \fdot{navy}{5}{0} \fdot{goldacc}{0}{5}
\end{tikzpicture}
\end{minipage}
\\
& \notesprompt{We work these together.}
\begin{probgrid}{|Y|Y|}
\pcell{11}{\textbf{Vertex form.} $y=-(x-2)^2+9$: vertex, axis, direction --- and is $9$ a
maximum or a minimum?}{1.8cm}{} &
\pcell{12}{\textbf{Standard form.} $y=-x^2+4x+5$: the $y$-intercept straight off, then the axis
from $x=-\frac{b}{2a}$ (work below).}{1.8cm}{} \\ \hline
\end{probgrid}
\\
& \begin{work}
  x &= -\frac{b}{2a} = -\frac{4}{2(-1)} = 2 \\
  y &= -(2)^2 + 4(2) + 5 = 9
\end{work}
\\
& \begin{probgrid}*{|Y|Y|}
\pcell{13}{\textbf{Intercept form.} $y=-(x+1)(x-5)$: the zeros, and the axis at their
midpoint.}{1.8cm}{} &
\pcell{14}{\textbf{The crux.} Do your three axes agree, and why must they? A classmate read
the zeros as $1$ and $-5$: what went wrong?}{1.8cm}{} \\ \hline
\end{probgrid}
\\ \hline
\end{guidednotes}

% ── THE NOTES END AT GUIDED PRACTICE ─────────────────────────────────────────
% There is deliberately no "you do" here: ../homework/ is the individual practice,
% started in class in the last ten minutes while the teacher circulates.

\end{document}
